\section*{अध्याय \ref{ch:g12:cells-and-electrolysis} --- विद्युत-रासायनिक सेल और विद्युत-अपघटन}

\begin{solution}{exo:g12:cells-and-electrolysis:1}
जस्ते का इलेक्ट्रोड \omterm{def:g12:cells-and-electrolysis:anode}{ऐनोड} है, \ce{Zn -> Zn^{2+} + 2e-}; ताँबे का
इलेक्ट्रोड \omterm{def:g12:cells-and-electrolysis:anode}{कैथोड} है, \ce{Cu^{2+} + 2e- -> Cu}।
\end{solution}

\begin{solution}{exo:g12:cells-and-electrolysis:2}
% ledger: const:F
$Q = 0.20 \times 1800 = \qty{360}{C}$; $n(e^-) = 360/96485 =
\qty{3.7e-3}{mol}$।
\end{solution}

\begin{solution}{exo:g12:cells-and-electrolysis:3}
वह परिपथ पूरा करता है और हर विलयन को विद्युत रूप से उदासीन रखता है, उसके
\omterm{def:g9:ions:ion}{आयन} कक्षों में जाते हैं। उसके बिना परिपथ खुला है: कोई
धारा नहीं बहती, वोल्टमीटर शून्य पढ़ता है और बल्ब अँधेरा रहता है।
\end{solution}

\begin{solution}{exo:g12:cells-and-electrolysis:4}
$2.5 \times 3600 = \qty{9000}{C}$।
\end{solution}

\begin{solution}{exo:g12:cells-and-electrolysis:5}
नहीं: अभिक्रिया स्वतः होने वाली दिशा के विपरीत दिशा में चलती है।
ऊर्जा जनित्र देता है।
\end{solution}

\begin{solution}{exo:g12:cells-and-electrolysis:6}
\omterm{def:g12:cells-and-electrolysis:anode}{ऐनोड} (ताँबा, $-$): \ce{Cu -> Cu^{2+} + 2e-}। \omterm{def:g12:cells-and-electrolysis:anode}{कैथोड} (चाँदी, $+$):
\ce{Ag+ + e- -> Ag}। समग्र: \ce{Cu + 2Ag+ -> Cu^{2+} + 2Ag}।
\end{solution}

\begin{solution}{exo:g12:cells-and-electrolysis:7}
तीन वर्ष $3 \times 365 \times 24 = \qty{26280}{h}$ हैं;
$0.010 \times 26280 = \qty{263}{mA.h}$।
\end{solution}

\begin{solution}{exo:g12:cells-and-electrolysis:8}
% ledger: const:F, aw:Ag
$Q = 0.50 \times 2400 = \qty{1200}{C}$; $n(e^-) = n(\ce{Ag}) =
1200/96485 = \qty{0.0124}{mol}$; $m = 0.0124 \times 107.9 =
\qty{1.34}{g}$।
\end{solution}

\begin{solution}{exo:g12:cells-and-electrolysis:9}
तार में धारा ताँबे ($+$) से जस्ते ($-$) की ओर जाती है,
\omterm{def:g9:inside-the-atom:nucleus}{इलेक्ट्रॉनों} के विपरीत। \ce{K+} \omterm{def:g9:ions:ion}{आयन} ताँबे वाले कक्ष की ओर जाते हैं,
जो \ce{Cu^{2+}} खोता है; \ce{NO3-} \omterm{def:g9:ions:ion}{आयन} जस्ते वाले
कक्ष की ओर, जो \ce{Zn^{2+}} पाता है।
\end{solution}

\begin{solution}{exo:g12:cells-and-electrolysis:10}
दोनों में \omterm{def:g12:cells-and-electrolysis:anode}{ऐनोड} वह है जहाँ \omterm{def:g11:redox:oxidation}{ऑक्सीकरण} होता है। सेल में \omterm{def:g12:cells-and-electrolysis:anode}{ऐनोड}
$-$ है, अभिक्रिया स्वतः होती है और रासायनिक ऊर्जा
विद्युत ऊर्जा बनती है; \omterm{def:g12:cells-and-electrolysis:electrolysis}{विद्युत-अपघटन} में \omterm{def:g12:cells-and-electrolysis:anode}{ऐनोड} जनित्र के $+$ से जुड़ा होता है,
अभिक्रिया कराई जाती है और विद्युत ऊर्जा
रासायनिक ऊर्जा बनती है।
\end{solution}

\begin{solution}{exo:g12:cells-and-electrolysis:11}
\qty{12}{mL}: समीकरण \ce{2H2O -> 2H2 + O2} डाइऑक्सीजन के एक \omterm{def:g7:atoms-and-molecules:molecule}{अणु} पर
डाइहाइड्रोजन के दो \omterm{def:g7:atoms-and-molecules:molecule}{अणु} देता है, और गैस की समान मात्राएँ समान
आयतन घेरती हैं।
\end{solution}

\begin{solution}{exo:g12:cells-and-electrolysis:12}
% ledger: aw:Zn, const:F
$n(\ce{Zn}) = 5.0/65.4 = \qty{0.076}{mol}$;
$n(\ce{Cu^{2+}}) = 0.100 \times 0.50 = \qty{0.050}{mol}$; वे एक-एक के
अनुपात में अभिक्रिया करते हैं, इसलिए कॉपर \omterm{def:g9:ions:ion}{आयन} पहले समाप्त होते हैं। $n(e^-) = 2 \times 0.050
= \qty{0.100}{mol}$, $Q = \qty{9.65e3}{C} = \qty{2.68}{A.h}$।
\qty{50}{mA} पर: $2.68/0.050 = \qty{54}{h}$।
\end{solution}

\begin{solution}{exo:g12:cells-and-electrolysis:13}
$Q = 2.0 \times 3600 = \qty{7200}{C}$; $E = 7200 \times 1.5 =
\qty{1.08e4}{J}$, अर्थात $1.5 \times 2.0 = \qty{3.0}{W.h}$।
\end{solution}

\begin{solution}{exo:g12:cells-and-electrolysis:14}
% ledger: const:F
$Q = \qty{7200}{C}$, $n(e^-) = \qty{0.0746}{mol}$। प्रति
\ce{H2} दो \omterm{def:g9:inside-the-atom:nucleus}{इलेक्ट्रॉन}: $n(\ce{H2}) = \qty{0.0373}{mol}$, $V = \qty{0.90}{L}$। प्रति
\ce{O2} चार: $n(\ce{O2}) = \qty{0.0187}{mol}$, $V = \qty{0.45}{L}$।
\end{solution}

\begin{solution}{exo:g12:cells-and-electrolysis:15}
\omterm{def:g12:cells-and-electrolysis:anode}{ऐनोड} पर घुला ताँबा जमा हुए ताँबे के बराबर है, क्योंकि दोनों जगह वही
\omterm{def:g9:inside-the-atom:nucleus}{इलेक्ट्रॉन} गुज़रते हैं: $2.84/3.00 = \qty{94.7}{\percent}$ ताँबा।
\end{solution}

\begin{solution}{pb:g12:cells-and-electrolysis:1}
% ledger: const:F, const:NA, aw:Li, dch:CH4
\textbf{1.} \omterm{def:g12:cells-and-electrolysis:anode}{ऐनोड}: वहाँ लिथियम ऑक्सीकृत होता है।

\textbf{2.} ग्रेफ़ाइट इलेक्ट्रोड, यानी \omterm{def:g12:cells-and-electrolysis:anode}{ऐनोड}।

\textbf{3.} \omterm{def:g9:inside-the-atom:nucleus}{इलेक्ट्रॉन} फ़ोन से होकर ग्रेफ़ाइट इलेक्ट्रोड से
ऑक्साइड इलेक्ट्रोड तक जाते हैं; लिथियम \omterm{def:g9:ions:ion}{आयन} विद्युत-अपघट्य को उसी
दिशा में, ग्रेफ़ाइट से ऑक्साइड की ओर, पार करते हैं।

\textbf{4.} चार्जर इसकी अभिक्रिया को उलटी दिशा में चला सकता है।

\textbf{5.} \qty{4.000}{A.h}, अर्थात $4.000 \times 3600 =
\qty{14400}{C}$।

\textbf{6.} $14400/96485 = \qty{0.1492}{mol}$।

\textbf{7.} $0.1492 \times \num{6.022e23} = \num{8.99e22}$ \omterm{def:g9:inside-the-atom:nucleus}{इलेक्ट्रॉन}।

\textbf{8.} $4.000/0.25 = \qty{16}{h}$।

\textbf{9.} $0.20 \times 14400 = \qty{2880}{C}$।

\textbf{10.} $E = 14400 \times 3.7 = \qty{5.3e4}{J}$, लगभग \qty{53}{kJ}।

\textbf{11.} $53280/3600 = \qty{14.8}{W.h}$।

\textbf{12.} \qty{1}{g} मेथेन \omterm{def:g4:what-a-fire-needs:burning}{जलाने} पर पूरी भरी फ़ोन बैटरी में संचित ऊर्जा से
थोड़ी अधिक ऊर्जा (\qty{55.7}{kJ}) मुक्त होती है।

\textbf{13.} $1000/14.8 = 67.6$: 67 पूरी चार्जिंग।

\textbf{14.} \ce{Li+ + e- -> Li}: एक \omterm{def:g11:redox:oxidation}{अपचयन}।

\textbf{15.} चार्जर अभिक्रिया को स्वतः होने वाली दिशा के विपरीत दिशा में
चलने पर मजबूर करता है।

\textbf{16.} $14400/2.0 = \qty{7200}{s}$, \qty{2.0}{h}।

\textbf{17.} प्रति \omterm{def:g9:inside-the-atom:nucleus}{इलेक्ट्रॉन} एक लिथियम \omterm{def:g9:ions:ion}{आयन}: \qty{0.1492}{mol}।

\textbf{18.} $0.1492 \times 6.9 = \qty{1.03}{g}$।

\textbf{19.} एक पूरी चार्जिंग में इलेक्ट्रोडों के बीच लगभग \qty{1.03}{g} लिथियम
आता-जाता है।
\end{solution}
